% =============================================================
% CHAPTER 6 — RESULTS AND DISCUSSION
% Purpose: present your findings (results) and interpret them (discussion). This
% is where you actually ANSWER the research questions. Keep "what the numbers
% are" and "what the numbers mean" clearly distinguishable -- either as separate
% sections or clearly separated paragraphs.
% Tip: lead with the headline result. Every table/figure must be referred to in
% the text and interpreted -- never drop a table without saying what to see in it.
% =============================================================
\chapter{Results and Discussion}
\label{chap:results}

\guide{Open with a short paragraph stating the headline finding before any
tables, so a hurried reader gets the message immediately.}

\section{Overall Detection Performance}
\guide{Present the main results table: models (rows) vs metrics (columns),
aggregated over datasets/prompts as appropriate. Then INTERPRET it: which models
do best, are any above/below chance, how big are the gaps. Refer to the table by
\Cref{tab:main-results}.}

\begin{table}[t]
	\centering
	\caption{Detection performance (\%) across models.}
	\label{tab:main-results}
	\begin{tabular}{lcccc}
		\toprule
		Model & Accuracy & Precision & Recall & F1 \\
		\midrule
		\dots & -- & -- & -- & -- \\
		\bottomrule
	\end{tabular}
\end{table}

\section{Effect of Prompting Strategy (RQ2)}
\guide{Show how performance changes across your prompt variants. A grouped table
or bar chart works well. Discuss which strategies helped, which hurt, and why you
think so. Tie the discussion explicitly back to RQ2.}

\section{Results by Dataset and Manipulation Type (RQ3)}
\guide{Break results down by dataset (FOWS vs GOTCHA) and by generator/manipulation
family. This reveals generalisation behaviour -- the core motivation. Highlight
where models systematically fail.}

\section{Error Analysis}
\guide{Go beyond aggregate numbers. Show representative failure cases (with
images if licensing allows), categorise common error types, and quantify
parse-failure vs genuine misclassification. This qualitative analysis is often
what distinguishes a good thesis from a results dump.}

\section{Discussion}
\guide{Step back and answer each research question in plain language, supported by
the evidence above. What do the results imply about using MLLMs as face-swap
detectors? Compare against the related work from \Cref{chap:relatedwork}. Be
balanced -- discuss both what works and what does not.}

\section{Limitations and Threats to Validity}
\guide{State honestly what limits the conclusions: dataset coverage, model
selection, prompt sensitivity, single-image (not video) scope, evaluation
choices. Acknowledging limitations strengthens, not weakens, a thesis.}
